\chapter{Phony Targets and Project Conventions}\label{ch:mk-phony}

Not every target is a file. \cmd{make clean}, \cmd{make test} and \cmd{make
	install} name \emph{actions}, and make's file-and-timestamp model needs to be
told so.

\section{\texorpdfstring{\texttt{.PHONY} and Why It Matters}{.PHONY and Why It Matters}}\label{sec:mk-phonyt}

\begin{makecode}
clean:
    rm -f app *.o
\end{makecode}

This works until the day someone creates a file called \file{clean}:

\begin{shell}
$ touch clean
$ make clean
make: 'clean' is up to date.
\end{shell}

Make did exactly what it promised. The target \texttt{clean} exists, it has no
prerequisites, so nothing is out of date and the recipe does not run. The fix
is to declare that this target is not a file:

\begin{makecode}
.PHONY: all clean install test

clean:
    rm -f app *.o
\end{makecode}

\begin{definition}[Phony target]
	A target listed as a prerequisite of the special target \texttt{.PHONY}. Make
	then never checks the filesystem for it: the recipe runs unconditionally, and
	no timestamp comparison is made.
\end{definition}

\begin{note}
	The stray file is the story everyone tells, and it is rare. The reason to use
	\texttt{.PHONY} every time is the second, mundane one: it is faster --- make
	skips the implicit-rule search entirely for phony targets --- and it documents
	intent. A reader can tell at a glance which of your targets produce files.
\end{note}

\begin{moral}
	Every target that does not name a file it creates goes in \texttt{.PHONY}.
	Declare it once, near the bottom, listing all of them.
\end{moral}

\section{The Standard Target Names}\label{sec:mk-standard}

These names are close to universal. Using them means anyone can build your
project without reading the Makefile.

\begin{keys}[0.20]
	\texttt{all}       & The default build. Should be the first rule (\autoref{sec:mk-goal}) \\
	\texttt{clean}     & Remove everything the build produced                     \\
	\texttt{distclean} & \dots{} and the configuration too, back to a fresh checkout \\
	\texttt{install}   & Copy the results into \texttt{\$(PREFIX)}                \\
	\texttt{uninstall} & Remove them again                                        \\
	\texttt{test}, \texttt{check} & Run the test suite                            \\
	\texttt{dist}      & Build a release tarball                                  \\
	\texttt{help}      & Print the available targets                              \\
\end{keys}

\begin{makecode}
PREFIX ?= /usr/local
BINDIR := $(PREFIX)/bin

install: app
    install -d $(DESTDIR)$(BINDIR)
    install -m 755 app $(DESTDIR)$(BINDIR)/app

uninstall:
    rm -f $(DESTDIR)$(BINDIR)/app

test: app
    ./run-tests.sh

clean:
    rm -f app $(OBJ) $(DEP)

.PHONY: all install uninstall test clean
\end{makecode}

\begin{note}
	\texttt{DESTDIR} is prepended to every installed path and is normally empty.
	Distribution packagers set it to a staging directory, so that
	\cmd{make install DESTDIR=/tmp/pkg} lays out a complete tree under
	\file{/tmp/pkg} without touching the real \file{/usr/local}. It costs nothing
	to support and packagers will thank you.
\end{note}

\begin{eg}[A self-documenting \texttt{help}]
	A pleasant trick: annotate targets with a marked comment and let
	\cmd{awk} (\autoref{sec:sh-awk}) extract them.
\begin{makecode}
help:  ## show this message
    @awk -F':.*##' '/^[a-z-]+:.*##/ {printf "  %-12s %s\n", $$1, $$2}' \
        $(MAKEFILE_LIST)

build: ## compile everything
test:  ## run the test suite
\end{makecode}
	Note \texttt{\$\$1}: make expands \texttt{\$} first, so the \cmd{awk} program
	needs its dollars doubled to survive --- exactly the hazard named in
	\autoref{sec:mk-notation}.
\end{eg}

\section{Order-Only Prerequisites and Directories}\label{sec:mk-orderonly}

\autoref{sec:mk-vpath} left a problem: \file{build/} must exist before anything
is written into it, but it must not trigger rebuilds.

The naive fix does not work:

\begin{makecode}
$(BUILD)/%.o: %.c $(BUILD)
    $(CC) -c -o $@ $<
\end{makecode}

A directory's timestamp changes every time a file is written into it. So the
directory becomes newer than the objects, and every object rebuilds on every
run.

\begin{definition}[Order-only prerequisite]
	A prerequisite written after a \texttt{|} in the prerequisite list. It must be
	built before the target, but its timestamp is ignored when deciding whether the
	target is out of date.
\end{definition}

\begin{makecode}
$(BUILD)/%.o: %.c | $(BUILD)
    $(CC) $(CPPFLAGS) $(CFLAGS) -c -o $@ $<

$(BUILD):
    mkdir -p $@
\end{makecode}

\begin{note}
	Read the \texttt{|} as ``and, before any of that, make sure this exists''.
	Creating output directories is the overwhelmingly common use; the other is a
	code generator that must run first but whose own timestamp should not force
	everything downstream to rebuild.
\end{note}

\section{Recursive Make, and Its Problems}\label{sec:mk-recursive}

A large project is often split into a Makefile per directory, each invoking make
in its subdirectories:

\begin{makecode}
SUBDIRS := src lib tests

all:
    for d in $(SUBDIRS); do $(MAKE) -C $$d || exit 1; done
\end{makecode}

\begin{note}
	\texttt{\$(MAKE)} rather than \cmd{make}: it carries the parent's flags through
	(including \cmd{-j} and \cmd{-n}) and lets make recognise the sub-build as its
	own. \cmd{-C dir} changes directory first. Writing \cmd{make} literally here
	breaks parallel builds and \cmd{make -n} alike.
\end{note}

This arrangement is common, and it is a known mistake. Peter Miller's 1997
paper \emph{Recursive Make Considered Harmful} \cite{miller1998recursive} sets
out why, and the argument has not weakened:

\begin{enumerate}[(i)]
	\item \textbf{No sub-make sees the whole graph.} Each knows only its own
	      directory, so a dependency crossing a directory boundary cannot be expressed
	      --- only guessed at by ordering the subdirectories by hand.
	\item \textbf{The order is a lie.} \texttt{src} before \texttt{tests} is a
	      hand-maintained approximation of a graph make could have computed exactly.
	\item \textbf{Parallelism is crippled.} \cmd{make -j8} cannot overlap two
	      subdirectories that a shell \cmd{for} loop runs in sequence, so the build is
	      as slow as its longest chain.
	\item \textbf{Builds can be wrong, not merely slow.} A missing cross-directory
	      edge means make declares success with a stale object linked in --- the worst
	      failure a build system has.
\end{enumerate}

\begin{moral}
	Prefer one Makefile that \cmd{include}s fragments from each directory over one
	Makefile per directory that calls make. Each fragment contributes its rules to
	a single graph, make sees all of it, and \cmd{-j} works properly.
\end{moral}

\begin{makecode}
# top-level Makefile: one graph, several files
include src/module.mk
include lib/module.mk
include tests/module.mk

all: $(ALL_TARGETS)
\end{makecode}

\begin{remark}
	This is also roughly what CMake and Ninja do: the description is written per
	directory, and the graph that gets executed is a single flat one. The idea did
	win --- it just stopped being spelled \cmd{make -C}.
\end{remark}
